\section{Setting and what is attainable when everything is fixed}\label{sec:theory}
\paragraph{Setting.}
Let $(X_i,Y_i)$, $i=1,\dots,m$, be i.i.d.\ with $m=n+N$. A uniformly random subset $L$ of size $n$ is labeled; for $i\notin L$ the label is missing (missing completely at random, MCAR, by design of the labeling). The pre-trained predictor $f$ is fixed and not trained on any of these data; write $Z_i=f(X_i)\in\mathbb R$ (the argument below only uses that $Z_i$ is a measurable function of $X_i$). The analyst sees $Z_i$ for all $m$ points and $Y_i$ for $i\in L$, and wants a $(1-\alpha)$ interval for $\theta=\mathbb E[Y]$. We chose the mean because it is the simplest target for which the efficiency question has a closed-form answer, it is the target for which PPI was first demonstrated, and it is what all nine example data sets in \texttt{ppi-python} are used for. Write
\[
g(z)=\mathbb E[Y\mid Z=z],\quad \sigma^2=\mathbb E\,\mathrm{Var}(Y\mid Z),\quad V=\mathrm{Var}(Z),\quad \beta=\tfrac{\mathrm{Cov}(Y,Z)}{V},\quad h(z)=g(z)-\mathbb E g-\beta (z-\mathbb E Z),
\]
so $h$ is the part of the calibration curve $g$ that is \emph{not} affine in the predictor output, and $v_h=\mathrm{Var}\,h(Z)$. Let $s=v_h/(\sigma^2+v_h)$ be the \emph{nonlinear share} of the residual variance left by the best affine recalibration. All variances below are asymptotic variances of the estimator of $\theta$ (i.e.\ squared interval half-widths up to the factor $z_{1-\alpha/2}^2$); $n,N\to\infty$ with $n/m\to\pi\in(0,1]$.

\paragraph{Existing estimators.}
PPI with weight $\lambda$ uses $\hat\theta_\lambda=\frac1N\sum_{j\notin L}\lambda Z_j+\frac1n\sum_{i\in L}(Y_i-\lambda Z_i)$ \citep{angelopoulos2023ppi,angelopoulos2023ppipp}; $\lambda=1$ is the original PPI, $\lambda=0$ the classical sample mean, and PPI++ picks $\lambda$ from the data to minimize variance. Its asymptotic variance is $\mathrm{Var}(Y-\lambda Z)/n+\lambda^2V/N$.

\begin{proposition}[Variance of the PPI family]\label{prop:ppi}
Minimizing over $\lambda\in\mathbb R$ gives $\lambda^\star=\beta N/m$ and
$V_{\mathrm{PPI++}}=\dfrac{\sigma^2+v_h}{n}+\dfrac{\beta^2V}{m}$. For fixed $\lambda$,
$V_\lambda-V_{\mathrm{PPI++}}=V\frac{m}{nN}(\lambda-\lambda^\star)^2$; in particular PPI with $\lambda=1$ is suboptimal by $V\frac{m}{nN}(1-\beta N/m)^2$, which vanishes only if $\beta N/m=1$.
\end{proposition}

\begin{proposition}[Efficiency bound with the predictor output as the only information]\label{prop:bound}
In the model where $(Z,Y)$ has an unrestricted distribution (finite variance), any regular asymptotically linear estimator of $\theta$ based on $\{Z_i\}_{i\le m}$ and $\{Y_i\}_{i\in L}$ has asymptotic variance at least
\[V_{\mathrm{eff}}=\frac{\sigma^2}{n}+\frac{\mathrm{Var}\,g(Z)}{m}=\frac{\sigma^2}{n}+\frac{\beta^2V+v_h}{m}.\]
It is attained by $\hat\theta=\frac1n\sum_{L}(Y_i-g(Z_i))+\frac1m\sum_{i\le m}g(Z_i)$, whose influence function is $\frac{R}{\pi}(Y-g(Z))+g(Z)-\theta$.
\end{proposition}
\begin{proof}[Proof sketch]
This is the standard semiparametric result for the mean under MCAR missingness with known labeling probability $\pi$ \citep{robins1994estimation,bickel1993efficient,tsiatis2006semiparametric}: with the nonparametric model for $(Z,Y)$, the observed-data influence function of $\theta$ is unique and equals $\frac{R}{\pi}(Y-\theta)-\frac{R-\pi}{\pi}\mathbb E[Y-\theta\mid Z]$. Its variance per observation is $\mathbb E(Y-g)^2/\pi+\mathrm{Var}\,g$, which in units of the sample sizes is the displayed expression. The case $Z=\text{const}$ recovers $\mathrm{Var}(Y)/n$, and $N=\infty$ recovers $\sigma^2/n$ \citep{zhang2019semi}.
\end{proof}

\begin{corollary}[When PPI++ cannot be improved]\label{cor:gap}
$V_{\mathrm{PPI++}}-V_{\mathrm{eff}}=v_h\,\dfrac{N}{n\,m}$, hence
\[
\frac{V_{\mathrm{eff}}}{V_{\mathrm{PPI++}}}=1-\frac{N}{m}\cdot\frac{s}{1+\kappa}\ \ \ge\ 1-\frac{N}{m}\,s,\qquad \kappa=\frac{n}{m}\frac{\beta^2V}{\sigma^2+v_h}\ge 0 .
\]
Thus (i) PPI++ is asymptotically efficient if and only if $g$ is a.s.\ affine in $Z$ ($v_h=0$) or $N/m\to0$ (essentially no unlabeled data); (ii) whenever the predictor output takes two values (a hard 0/1 prediction), $g$ is automatically affine, so no estimator that uses only $Z$ can beat PPI++; (iii) the interval-width reduction is at most $1-\sqrt{1-sN/m}\le 1-\sqrt{1-s}$, e.g.\ about $5\%$ for $s=0.10$ and about $1\%$ for $s=0.02$.
\end{corollary}
The corollary follows by subtracting the two expressions in Props.~\ref{prop:ppi} and~\ref{prop:bound}. It says exactly \emph{what} room remains below PPI++ with labeled set, unlabeled set and predictor all fixed: only the non-affine part of the calibration curve $\mathbb E[Y\mid f]$, discounted by the fraction $N/m$ of unlabeled data. It also shows that the large-$N$ regime that motivates PPI is where the headroom is largest. Using side features $X$ beyond $Z=f(X)$ lowers the bound further to $\mathbb E\,\mathrm{Var}(Y|X)/n+\mathrm{Var}(\mathbb E[Y|X])/m\le V_{\mathrm{eff}}$, but attaining it requires learning $\mathbb E[Y|X]$ from the labels, i.e.\ in effect replacing the predictor; we only quantify how large this extra gap is on the data at hand (\S\ref{sec:exp-x}).

\begin{remark}[Invariance]
A monotone reparametrization $Z\mapsto T(Z)$ (e.g.\ logits instead of probabilities) leaves $g$-based estimators with monotone or nonparametric $g$ unchanged, but changes $\beta$, $v_h$ and hence the width of PPI/PPI++. A predictor whose scores are not on the scale of $Y$ is therefore exactly the case with $v_h>0$.
\end{remark}
